\documentclass[11pt,a4paper]{article}

\usepackage[T1]{fontenc}
\usepackage[utf8]{inputenc}
\usepackage[english]{babel}
\usepackage{lmodern}
\usepackage{amsmath,amssymb,amsthm}
\usepackage[a4paper,margin=2.55cm]{geometry}
\usepackage{microtype}
\usepackage{booktabs,array}
\usepackage{fancyvrb}
\usepackage[hidelinks]{hyperref}
\usepackage{enumitem}
\usepackage{xurl}
\setlist[itemize]{leftmargin=*,topsep=3pt,itemsep=2pt}
\setlength{\parskip}{4pt}
\setlength{\parindent}{0pt}
\emergencystretch=2em

\newcommand{\Truth}{\operatorname{Truth}}
\newcommand{\Proof}{\operatorname{Proof}}
\newcommand{\ExactRealization}{\operatorname{ExactRealization}}
\newcommand{\FutureContract}{\operatorname{FutureContract}}
\newcommand{\FutureEquivalent}{\operatorname{FutureEquivalent}}
\newcommand{\Outcome}{\operatorname{Outcome}}
\newcommand{\Empty}{\mathsf{Empty}}
% Long technical identifiers must remain breakable so that hashes and claim
% names do not extend beyond the text area.
\newcommand{\code}[1]{\nolinkurl{#1}}
\newcommand{\Source}[2]{\href{https://github.com/JohnDoe-collab-stack/relational-perimeter/blob/1e5aeda745559d62a3819bb3213da7fce3c31ad1/#1}{\nolinkurl{#2}}}
\newcommand{\SourceAt}[3]{\href{https://github.com/JohnDoe-collab-stack/relational-perimeter/blob/1e5aeda745559d62a3819bb3213da7fce3c31ad1/#1\#L#2}{\nolinkurl{#3}}}
\newcommand{\MainAxis}{%
  \begin{gathered}
  \text{certified execution}\\
  \downarrow\\
  \bigl(C(e),\Gamma_e,\sigma_e\bigr)\\
  \downarrow\\
  \text{universal obligation}\\
  \text{over every exact realization of }C(e)
  \end{gathered}}
\newcommand{\FormalRevision}{1e5aeda745559d62a3819bb3213da7fce3c31ad1}

\newtheorem{theorem}{Theorem}[section]
\newtheorem{proposition}[theorem]{Proposition}
\newtheorem{lemma}[theorem]{Lemma}
\theoremstyle{remark}
\newtheorem{remark}[theorem]{Remark}

\title{From Gödelian Execution to Universal Obligations\\
  of Exact Realization\\[5pt]
  \large Operational Constitution of Contracts, Separation, and Finite Progression}
\author{Formalization \code{relational-perimeter}}
\date{Formal revision \code{1e5aeda74555}\\
  \small Revised document version of October 23, 2026}

\begin{document}
\maketitle
\vspace{-1.4em}

\begin{abstract}
Within the broader formal theory of future contracts, operational grouping,
finite separation, directed readout preservation, and exact realization, the
result presented here is the specifically Gödelian construction that derives
a universal obligation from a certified execution and renews this
per-execution construction through a finite progression of proof regimes.

We present a result whose object is not only the unprovability produced by
Gödelian diagonalization, but the obligation that its \emph{certified
execution} imposes on every exact realization of a contract constituted
from that same execution. At each certified arithmetic stage, an executed
production supplies a sentence $G$ and witnesses of $\Truth(G)$ and
$\Proof_T(G)\to\Empty$. The runtime constituted from that execution computes
two terminals. The same execution canonically derives its contract, the
package containing the two runs and their operational collision, and then a
finite separator between the terminals thus produced. It follows that, for
every memory type and every exact realization of the contract, their memory
images are distinct. This separation is not a field of the definition of
exactness: it is derived from the certified endpoint witnesses, the terminal
readouts, and their preservation laws. In parallel, an
operational grouping with a common terminal memory is effectively
constructed when global readout exactness is not required, while an exact
identity realization ensures that the result is non-vacuous. Finally, a
finite Gödelian progression retains every execution, extends the history
exactly by the produced sentence, and then renews the contract and its
universal obligation at the new stage. The development therefore constructs,
for every finite natural height, a chain in which each retained execution
generates a universal law governing all its admissible exact realizations.
Its exact object is the per-execution derivation of a contract and its
obligation, together with their finite renewal along the executions that were
actually produced.
\end{abstract}

\noindent\textbf{Central object.}
\[
\boxed{\MainAxis}
\]
\noindent\textbf{Reference formalization.}
\href{https://github.com/JohnDoe-collab-stack/relational-perimeter}{\code{relational-perimeter}}\\
Revision \code{1e5aeda745559d62a3819bb3213da7fce3c31ad1}.

\noindent\textbf{Tracked claims.}
\begin{itemize}[leftmargin=1.25em,topsep=2pt,itemsep=0pt]
\item \code{GODEL_OPERATIONAL_RUNTIME_CLOSURE_OBSTRUCTION}
\item \code{GODEL_CONSTITUTIVE_THEORY_PROGRESSION}
\end{itemize}

\begin{center}
\fbox{\begin{minipage}{0.91\linewidth}
\textbf{Contract constituted by an execution $e$.}
\[
C(e)=\bigl(S_e,I,E,O,\operatorname{next},\operatorname{event},
  \operatorname{read},\operatorname{Allow},\operatorname{decision}\bigr),
\]
with $O=\mathrm{Type}$ and, for the two positions obtained from the same
produced Gödelian witness,
\[
\operatorname{read}_{C(e)}(t^{\mathrm v})=\Truth(G),\qquad
\operatorname{read}_{C(e)}(t^{\mathrm p})=\Proof_T(G).
\]
The contract is fixed before any choice of memory type or realization.
\end{minipage}}
\end{center}

\section{A change in the object of proof}

The production of a Gödelian sentence $G$ certified as true and unprovable
in a theory $T$ remains an indispensable ingredient. The result presented
here, however, does not stop at these two properties. The execution that
produces the witnesses serves to \emph{constitute a contract} $C$; the
contract fixes transitions, events, authorizations, and readouts; the
distinction certified by the witness then becomes a necessary property of
every realization that preserves this contract exactly.

We thereby move from a statement about the proof regime of $T$ to a
universal statement about the admissible exact representations of $C$.
Universality does not mean that all arbitrary representations are
admissible: the quantification ranges \emph{over those that realize the
contract of the execution exactly}.

\paragraph{Governing statement.}
At every finite stage $n$ of the progression under consideration, let
$e_n$ be the retained Gödelian execution, $C_n$ its contract,
$t_n^{\mathrm v}$ and $t_n^{\mathrm p}$ its truth and proof terminals, and
$q_A$ the memory projection of a realization $A$. Then:
\begin{equation}\label{eq:main}
\boxed{\begin{gathered}
C_n=\operatorname{godelClosureContract}(e_n),\\[3pt]
\forall M,\;\forall A:\ExactRealization(C_n,M),\\[2pt]
q_A(t_n^{\mathrm v})\neq q_A(t_n^{\mathrm p}).
\end{gathered}}
\end{equation}
The difference between the memories is a \emph{conclusion}, not a
constitutive clause of $\ExactRealization$.

\section{Constitution of the contract from the execution}
\label{sec:contract}

\subsection{Executed data and witnesses}

Let $T$ be a stage equipped with the certification required by the formal
interface. The executed Gödelian production $e$ contains a sentence $G$
and inhabitants of the types
\begin{equation}\label{eq:witnesses}
\Truth(G),\qquad \Proof_T(G)\longrightarrow\Empty.
\end{equation}
Here $\Proof_T(G)$ denotes the \emph{proof type} relative to the regime
$T$; the second term is not a proof of $G$ in $T$, but a function that
would reject any inhabitant of such a type.

The execution also serves as a parameter in the constitution of the two
poles of the runtime, its sources, and its authorized closure runs. These
runs compute the terminals $t^{\mathrm v}$ and $t^{\mathrm p}$, to which
the readouts
\begin{equation}\label{eq:readouts}
\operatorname{read}_C(t^{\mathrm v})=\Truth(G),\qquad
\operatorname{read}_C(t^{\mathrm p})=\Proof_T(G)
\end{equation}
are associated. These equalities concern values of the observation type,
which here is a Lean universe of types. They do not directly compare
Boolean truth values.

\subsection{The operational contract fixed before any realization}

In \SourceAt{RelationalPerimeter/Constitution/Continuation/Behavior.lean}{9-L14}{Behavior.lean}
(lines 9--14), a \code{FutureContract} contains, for types $S$ (sources),
$I$ (requests), $E$ (events), and $O$ (observations), the operations
\[
\begin{aligned}
\operatorname{next}&:S\times I\to S, &
\operatorname{event}&:S\times I\to E,\\
\operatorname{read}&:S\to O, &
\operatorname{Allow}&:S\times I\to\mathrm{Type},
\end{aligned}
\]
together with an effective authorization-or-rejection decision for each
pair $(s,i)$. The contract specifies observations even when a request is
rejected.

For the Gödelian execution, the definition \code{godelClosureContract e}
in \SourceAt{RelationalPerimeter/Constitution/Godel/OperationalClosure.lean}{150-L157}{OperationalClosure.lean}
(lines 150--157) fixes the five components of the contract: \code{next},
\code{event}, \code{read}, \code{Allow}, and \code{decision}. Its readout
is determined by the pole of the source: \code{Truth(G)} for the left pole
and \code{Proof(G)} for the right pole. Neither a reduced memory type nor an
exact realization enters into the definition of this contract.

\subsection{Exact meaning of realization}

The structure \code{ExactRealization}, defined in
\SourceAt{RelationalPerimeter/Constitution/Continuation/Minimality.lean}{9-L23}{Minimality.lean}
(lines 9--23), supplies a reduced contract on a memory type $M$ and a
projection $q_A:S\to M$. It imposes bilateral preservation of
authorization proofs (with round-trip laws), commutation with transitions,
equality of events, and exactness of readouts. In particular:
\begin{equation}\label{eq:exactness}
\begin{aligned}
q_A(\operatorname{next}_C(s,i))
  &=\operatorname{next}_{A}(q_A(s),i),\\
\operatorname{event}_C(s,i)
  &=\operatorname{event}_{A}(q_A(s),i),\\
\operatorname{read}_C(s)
  &=\operatorname{read}_{A}(q_A(s)).
\end{aligned}
\end{equation}
This structure contains no field of the form
$q_A(t^{\mathrm v})\neq q_A(t^{\mathrm p})$. Preservation is fixed
generically; the necessity of preserving \emph{this particular
distinction} comes from the Gödelian execution.

\subsection{The canonically derived constitutive package}

It is useful to gather the data produced by one execution without turning
them into a new primitive. From every certified Gödelian execution $e$, the
following are canonically derived:
\begin{equation}\label{eq:constitutive-package}
e\longmapsto\bigl(C(e),\Gamma_e,\sigma_e\bigr).
\end{equation}
This notation summarizes a chain of constructions that remain separate in
the formalization; it does not denote an additional Lean structure. The
package $\Gamma_e$ retains the two nonempty authorized runs, their events,
and their actually computed terminals:
\[
\Gamma_e=
\bigl(r_e^{\mathrm v},r_e^{\mathrm p},
      t_e^{\mathrm v},t_e^{\mathrm p},\kappa_e\bigr),
\]
where $\kappa_e$ certifies the collision of the terminal projections in the
operational realization without global readout exactness. This collision
identifies neither the sources nor their constitutive profiles.

The final component is constructed on the same terminals:
\begin{equation}\label{eq:sigma-package}
\sigma_e:
\mathrm{FutureSeparator}
  \bigl(C(e),t_e^{\mathrm v},t_e^{\mathrm p}\bigr),
\qquad \sigma_e.\mathrm{requests}=[].
\end{equation}
It depends on the witnesses and runs of $e$, but on no choice of memory
type $M$ or exact realization $A$. The contract therefore fixes the
preservation rules, while $\sigma_e$ determines the particular distinction
that these rules oblige every exact realization to preserve. We obtain:
\begin{equation}\label{eq:package-obligation}
\boxed{
\sigma_e\Longrightarrow
\forall M,\ \forall A:\ExactRealization(C(e),M),\quad
q_A(t_e^{\mathrm v})\neq q_A(t_e^{\mathrm p}).}
\end{equation}
The order of dependencies is thus explicit: the execution constructs the
runs and their terminals; the executed witnesses certify $\sigma_e$;
exactness then makes separation mandatory for an arbitrary realization of
the already fixed contract.

\subsection{Two distinct consequences of the same contract}

The constituted contract supports two formal paths that must not be
confused. The first imposes a universal obligation on every exact
realization. The second starts from a positive operational grouping and
locates the incompatibility of an additional directed-preservation
requirement.

\begin{center}
\renewcommand{\arraystretch}{1.25}
\begin{tabular}{>{\centering\arraybackslash}p{0.44\linewidth}
                >{\centering\arraybackslash}p{0.44\linewidth}}
\textbf{Path A: universal separation}
& \textbf{Path B: collision and directed transport} \\
package $(C(e),\Gamma_e,\sigma_e)$ & package $(C(e),\Gamma_e)$ \\
$\downarrow$ & $\downarrow$ \\
finite separator $\sigma_e$ with request $[]$ & operational collision $\kappa_e$ \\
$\downarrow$ & $\downarrow$ \\
$\forall A:\ExactRealization(C(e),M)$,
$q_A(t^{\mathrm v})\neq q_A(t^{\mathrm p})$
& memory collision and additional directed preservation \\
& $\downarrow$ \\
& directed transport $\Truth(G)\to\Proof_T(G)$ \\
& $\downarrow$ using \code{trueWitness} and \code{unprovable} \\
& $\Empty$
\end{tabular}
\end{center}

In Path B, directed preservation is not a field of $C(e)$: it is the
additional requirement whose compatibility with the computed collision is
being tested.

\section{From the Gödelian witness to the finite separator}
\label{sec:separator}

\begin{lemma}[Incompatibility of the readouts]\label{lem:readout}
The witnesses in \eqref{eq:witnesses} establish, as an inequality between
the types used by the readout,
\[
\Truth(G)\neq\Proof_T(G).
\]
\end{lemma}
\begin{proof}
Assume an equality between these types. Transporting the witness of
$\Truth(G)$ along the equality gives an inhabitant of $\Proof_T(G)$.
Applying the witness $\Proof_T(G)\to\Empty$ produces an inhabitant of
$\Empty$, a contradiction. This is the form of
\code{godelWitnessReadouts\_ne},
\SourceAt{RelationalPerimeter/Constitution/Godel/OperationalClosure.lean}{636-L641}{OperationalClosure.lean},
lines 636--641.
\end{proof}

\paragraph{Distinguishing two request executions.}
The terminals are \emph{first computed} by authorized closure runs, each
carrying the nonempty request \code{[close]}. Once these terminals have
been obtained, the separator observes them with the empty request list
\code{[]}. This second list is not an empty execution that would have
produced the terminals: it observes readouts of terminals that have already
been produced.

The contract has a function $\operatorname{outcome}_C(s,w)$ assembling the
readout, authorizations, and events along the list $w$. By definition,
\[
\operatorname{outcome}_C(s,[]) =
  \operatorname{stop}(\operatorname{read}_C(s)).
\]
The preceding lemma and \eqref{eq:readouts} therefore give:
\begin{equation}\label{eq:separator}
\operatorname{outcome}_{C}(t^{\mathrm v},[])
\neq\operatorname{outcome}_{C}(t^{\mathrm p},[]).
\end{equation}
The object \code{godelTerminalSeparator e} constructs exactly this
\code{FutureSeparator}, with the field \code{requests := []}
(\SourceAt{RelationalPerimeter/Constitution/Godel/OperationalClosure.lean}{642-L667}{OperationalClosure.lean},
lines 642--667).

\section{Universal obligation of exact realization}
\label{sec:theorem}

\begin{lemma}[Generic preservation principle]\label{lem:generic}
For a contract $C$, an exact realization $A$ of $C$ on a memory type $M$,
and two sources $s,t$, one has
\[
q_A(s)=q_A(t)\quad\Longrightarrow\quad
\FutureEquivalent(C,s,t).
\]
Thus every finite separator of $s$ and $t$ rules out $q_A(s)=q_A(t)$.
\end{lemma}
\begin{proof}
Preservation of transition, authorization, event, and readout gives, by
induction on request lists, equality between the outcomes of the contract
and those of the reduced contract. Two equal projected memories then have
the same reduced outcomes for every list; by transport, the source outcomes
are equal. A separator supplies a list for which this equality is
impossible. These steps are the theorems
\code{ExactRealization.outcome\_exact},
\code{equal\_memory\_same\_futures}, and then
\code{distinct\_memory\_of\_separator} in
\SourceAt{RelationalPerimeter/Constitution/Continuation/Minimality.lean}{56-L75}{Minimality.lean}
and \SourceAt{RelationalPerimeter/Constitution/Continuation/ExactSeparation.lean}{20-L27}{ExactSeparation.lean}.
\end{proof}

\begin{theorem}[Universal separation generated by the Gödelian execution]
\label{thm:universal}
Let $e$ be a certified Gödelian execution and let $C(e)$ be the contract
\code{godelClosureContract e} that it constitutes. Let
$t^{\mathrm v}(e)$ and $t^{\mathrm p}(e)$ be the terminals actually
computed by the operational grouping of this execution. Then
\begin{equation}\label{eq:universal}
\boxed{\begin{gathered}
\forall M,\quad \forall A:\ExactRealization(C(e),M),\\
q_A(t^{\mathrm v}(e))\neq q_A(t^{\mathrm p}(e)).
\end{gathered}}
\end{equation}
\end{theorem}
\begin{proof}
The execution produces the witnesses in \eqref{eq:witnesses}, and its two
runs compute the terminals. Lemma~\ref{lem:readout} supplies the difference
between the readouts; the empty request supplies the finite separator
\eqref{eq:separator}. For an arbitrary exact realization $A$,
Lemma~\ref{lem:generic} rules out equality between their memory projections.
No choice of $M$ or $A$ enters into the construction of either the contract
or the separator. The statement therefore follows for every admissible
exact realization. In Lean, the theorem
\begin{center}
\code{godelEveryExactRealizationSeparatesTerminals}
\end{center}
is obtained by directly applying
\code{realization.distinct\_memory\_of\_separator} to
\code{godelTerminalSeparator execution}
(\SourceAt{RelationalPerimeter/Constitution/Godel/OperationalClosure.lean}{669-L676}{OperationalClosure.lean},
lines 669--676).
\end{proof}

\begin{remark}[Nature of universality]
The contract $C(e)$ \emph{depends} on the execution $e$. What is independent
of a particular realization is neither the existence of this execution nor
the choice of its contract, but the proof of the constraint
\eqref{eq:universal}. The quantification ranges over exact realizations of
$C(e)$, not over all possible models of Gödel.
\end{remark}

\begin{remark}[Minimal strength of the obstruction]
For these terminals, the separating list is empty: disagreement is
immediately visible in the readouts. The particular argument can therefore
be read directly through \code{read\_exact}. The generic future-preservation
lemma is more general than what is required for this separator alone. This
observation changes neither the constitution of the contract nor the
universal scope of the theorem.
\end{remark}

\section{Exact boundary: grouping, preservation, existence}
\label{sec:frontier}

The preceding theorem does not say that no grouping of the two branches can
be executed. The code precisely separates three statements.

\begin{proposition}[Operational grouping without global exactness]
An operational projection preserves transitions, events, and authorizations
while computing a common terminal memory after the two closure runs. It does
not assume global exactness of the readouts.
\end{proposition}
\begin{proof}[Construction]
The operational memory contains active states retaining the pole and
\code{grouped} states that no longer retain it. The closure transition
produces a common \code{grouped} state. The operational laws are established
in \code{godelOperationalRealizationWithoutRead};
\code{godelOperationalGroupingWithoutRead} executes the two runs and proves
their computed collision
(\SourceAt{RelationalPerimeter/Constitution/Godel/OperationalClosure.lean}{589-L602}{OperationalClosure.lean},
lines 165--272 and 358--372, 589--602). The reduced contract here reads a
constant \code{PUnit} observation: it does not claim to reproduce the two
source readouts.
\end{proof}

\begin{proposition}[Rejection of directed preservation in the presence of collision]
Adding to this grouping the directed preservation that converts the truth
readout into the proof readout would produce
$\Truth(G)\to\Proof_T(G)$ and then $\Empty$.
\end{proposition}
\begin{proof}
The grouping computes equality between the terminal memories. A
\begin{center}
\code{DirectedTerminalReadoutPreservation}
\end{center}
would then transport an inhabitant of the truth readout to the proof
readout. Applied to the true witness and followed by the unprovability
certificate, this conversion produces $\Empty$. This is the chain
\begin{center}
\code{godelComputedGroupingToProducedCompleteness}\\[2pt]
$\downarrow$\\[2pt]
\code{rejectGodelComputedGroupingToEmpty}
\end{center}
It is distinct from the proof of Theorem~\ref{thm:universal}.
\end{proof}

\begin{proposition}[Non-vacuity of exactness]
For every contract $C$ of type \code{FutureContract}, there is an exact
identity realization whose memory type is the source type. In particular,
the Gödelian contract has an exact realization.
\end{proposition}
\begin{proof}[Construction]
The definition \code{identityRealization} sets
$M=S$, $q_A=\mathrm{id}$, and takes the contract itself as the reduced
contract. The exactness laws reduce to reflexive equalities; see
\SourceAt{RelationalPerimeter/Constitution/Continuation/CommonMemoryGrouping.lean}{10-L19}{CommonMemoryGrouping.lean},
lines 10--19. The specialization is \code{godelIdentityRealization}, then
\code{heytingArithmeticIdentityRuntime} for the HA instance.
\end{proof}

These three results fix the boundary: \emph{operational fusion is
constructible without global readout exactness; fusion of the terminals is
forbidden in every exact realization of the contract; at least one such
exact realization effectively exists}. Non-vacuity is global: it does not
imply that every arbitrary memory type admits an exact realization.

\section{Constitutive renewal at every finite stage}
\label{sec:progression}

A \code{SoundGodelTheoryStage} retains a history $\mathcal H_n$, the
arithmetic regime $T_n$ that it presents, its certification, a master
Gödelian execution, and a proof that this execution is identical to the
canonical production of the regime. The procedure
\code{executeGodelTheoryProgressionFrom} consumes this retained execution.
It adds \emph{the exact occurrence} of the sentence that the execution
produced:
\begin{equation}\label{eq:successor}
\mathcal H_{n+1}
  =\operatorname{extend}\bigl(\mathcal H_n,
    \operatorname{occurrence}(G_n)\bigr).
\end{equation}
This notation denotes a constitutive extension of the history by the
occurrence actually produced. It denotes neither a primitive set-theoretic
union nor the addition of a sentence reconstructed independently. The
certification uses the truth witness of the source production. The newly
added occurrence supplies a verified proof of $G_n$ in $T_{n+1}$. This
proof is constructed \emph{by the added occurrence}: the Lean function
named \code{oldTruthToTargetProof} ignores its truth argument and returns
the already constructed target proof. It is therefore not a general
procedure that transforms truth into proof.

The same step executes a \emph{new} production $G_{n+1}$ relative to the
extended regime. The following data are then available:
\[
\begin{gathered}
\Proof_{T_n}(G_n)\to\Empty,
\qquad \Proof_{T_{n+1}}(G_n),\\
\Truth(G_{n+1}),\qquad
\Proof_{T_{n+1}}(G_{n+1})\to\Empty,\\
G_n\neq G_{n+1}.
\end{gathered}
\]
The inequality between the sentences is proved: identifying them would
transport the constructed proof of $G_n$ to the new unprovable sentence
$G_{n+1}$
(\SourceAt{RelationalPerimeter/Instances/HeytingArithmetic/GodelTheoryProgression.lean}{175-L207}{GodelTheoryProgression.lean},
lines 175--207).

\begin{theorem}[Finite renewal of the obligation]\label{thm:finite}
For every certified stage $S_0$ and every natural height $n$, the stage
$S_n=\operatorname{advance}(S_0,n)$ is constituted by finite iteration.
Let $e_n$ be the execution that it retains, let $C_n=C(e_n)$ be its
contract, and let $t_n^{\mathrm v}$ and $t_n^{\mathrm p}$ be its computed
terminals. The conclusion of \eqref{eq:main} holds at this stage. In
particular, this conclusion holds for the progression from the initial HA
stage.
\end{theorem}
\begin{proof}
The construction \code{advance} is recursive in $n$; the construction
\code{executedTrace} additionally retains every step, its target execution,
and its history. Its successor case consumes the terminal actually produced
by the preceding step instead of reconstructing an equivalent history from
a mere index. The universal property
\code{everyExactRealizationSeparatesCurrentTerminals} applies to every
\code{SoundGodelTheoryStage}, hence to the stage obtained at height $n$
(\SourceAt{RelationalPerimeter/Instances/HeytingArithmetic/GodelTheoryProgression.lean}{291-L448}{GodelTheoryProgression.lean},
lines 291--448).
\end{proof}

Extending the theory and realizing the contract remain two distinct
operations. No common exact realization is constructed between $C_n$ and
$C_{n+1}$; proof boundaries and obligations are reevaluated \emph{relative
to the contract proper to each stage}.

The executed trace retains the stages and the executions actually produced.
Universal obligations are not stored as fields of either the history or the
trace. They are derivable at every retained stage by applying
\code{everyExactRealizationSeparatesCurrentTerminals} to the execution proper
to that stage. Thus, at every stage of the finite trace, a renewed obligation
relative to its own contract is derivable.

\section{Contribution and exact scope of the result}

The formally presented contribution is the complete constitutive chain
\[
\boxed{\begin{gathered}
e\longmapsto\bigl(C(e),\Gamma_e,\sigma_e\bigr)\\[3pt]
\downarrow\\[3pt]
\forall M,\;\forall A:\ExactRealization(C(e),M),\\[2pt]
q_A(t_e^{\mathrm v})\neq q_A(t_e^{\mathrm p}).
\end{gathered}}
\]
It is realized by the Gödelian production, its finite separator, the universal
separation theorem, and their renewal at every finite height.

\paragraph{The discovery presented.}
A certified Gödelian execution supplies not only a sentence and properties
relative to its provability. From the same executed witness, it constitutes
its future contract, produces two runs and their terminals, constructs an
operational collision under the regime that does not impose global readout
exactness, and then constructs a finite separator on these same terminals.
This particular datum generates a consequence independent of every later
choice of memory: every admissible exact realization of the contract must
keep the projections of the two terminals distinct. The progression retains
the produced executions and renews this generation of a universal
obligation at every finite natural stage.

The decisive passage is not merely the final inequality between two
memories. From a single execution, the contract, the runs, the terminals,
and the separator are constructed, and the generic law of exactness is then
applied to them. The particular distinction imposed on realizations is
therefore neither a field of \code{ExactRealization} nor a separation chosen
before execution. It is generated by the objects actually produced and then
becomes universal relative to the contract thus constituted. The result
therefore turns the certified Gödelian witnesses into a representational
obligation governing every admissible exact realization.

The contribution therefore lies neither in the abstract lemma alone, which
states that a preserved observation must remain distinguishable, nor in the
production of a Gödelian sentence alone. It lies in their complete
constitutive articulation: the data to which the universal obligation
applies are precisely those produced and retained by the Gödelian execution
that constitutes the contract.

The result can be summarized in one sentence:
\begin{quote}
A certified Gödelian execution canonically derives a contract, two runs and
their terminals, an operational collision, and a finite separator of these
same terminals. Without postulating an inequality between their projected
memories, this separator imposes their separation in every admissible exact
realization; the progression of theories renews this obligation at every
constituted finite stage.
\end{quote}

\paragraph{Exact scope.}
The theorem concerns the contracts constituted by the certified executions
of the formalization and all their admissible exact realizations. The
progression is constructed for every finite natural height; it retains the
occurrences and executions actually produced and then reevaluates the
obligation relative to the contract proper to the new stage. Its object does
not include a limit stage, a transfinite union, a final complete theory, or
a transversal realization between successive contracts. The proved
separation concerns certified unprovability of the produced sentence; it
does not replace a separate proof of irrefutability.

\paragraph{Scientific positioning.}
Executed diagonalization, observational separation, contracts independent
of their implementations, and progressions of theories each have
established antecedents. The contribution claimed here is none of these
ingredients in isolation. In the comparative primary corpus examined, no
antecedent was identified that constructs, from the same executed Gödelian
witness, the complete chain
\[
e\longmapsto\bigl(C(e),\Gamma_e,\sigma_e\bigr)
\longmapsto
\forall M\,\forall A:\ExactRealization(C(e),M),\quad
q_A(t_e^{\mathrm v})\neq q_A(t_e^{\mathrm p}),
\]
while simultaneously establishing the possibility of operational grouping
without global exactness, the non-vacuity of the quantification through a
positive exact realization, and then renewing this obligation on the finite
trace actually produced. This positioning locates the contribution; it
replaces neither the formal chain nor the execution-first axis of exposition.

\section{Lean core and dependency traceability}
\label{sec:lean}

The central formula gathers four distinct formal constructions:
\begin{center}
\renewcommand{\arraystretch}{1.2}
\begin{tabular}{>{\raggedright\arraybackslash}p{0.28\linewidth}
                >{\raggedright\arraybackslash}p{0.61\linewidth}}
\toprule
\textbf{Manuscript notation} & \textbf{Lean construction} \\
\midrule
$C(e)$ & \code{godelClosureContract execution} \\
$\Gamma_e$ & \code{godelOperationalGroupingWithoutRead execution} \\
$\sigma_e$ & \code{godelTerminalSeparator execution} \\
universal obligation & \code{godelEveryExactRealizationSeparatesTerminals} \\
\bottomrule
\end{tabular}
\end{center}

The triple $(C(e),\Gamma_e,\sigma_e)$ is manuscript notation for this
composition of existing constructions. It does not denote an additional
Lean structure.

The passage that makes the change of object visible is the theorem
\begin{center}
\code{godelEveryExactRealizationSeparatesTerminals}
\end{center}
in
\SourceAt{RelationalPerimeter/Constitution/Godel/OperationalClosure.lean}{669-L676}{OperationalClosure.lean}
(lines 669--676). Its proof term is exactly:
\begin{Verbatim}[fontsize=\small]
realization.distinct_memory_of_separator (godelTerminalSeparator execution)
\end{Verbatim}
The progression reaches the following theorem, valid for an arbitrary
certified stage:
\begin{center}
\code{everyExactRealizationSeparatesCurrentTerminals}
\end{center}
Its \code{realization} hypothesis is quantified over every exact realization
of the contract of the execution retained by the stage.

\paragraph{Decisive dependencies, in their logical order.}
\begin{enumerate}[leftmargin=*,itemsep=3pt]
\item \textbf{Constitution.} \code{godelClosureContract} is defined from the execution, without a chosen memory.
\item \textbf{Production of $\Gamma_e$.} \code{godelOperationalGroupingWithoutRead} executes the two nonempty runs, retains their events and terminals, and then constructs their operational collision without global readout exactness.
\item \textbf{Construction of $\sigma_e$.} \code{godelTerminalSeparator} reuses the terminals of these runs, fixes the request \code{[]}, and proves disagreement between their outcomes from the witnesses of the same execution.
\item \textbf{Passage to the general result.} \code{distinct_memory_of_separator} turns the separator into a constraint on every exact realization, without choosing either the memory or the realization when $\sigma_e$ is constructed.
\item \textbf{Boundary.} The collision in $\Gamma_e$ is positive in the operational realization without exact readout, while \code{identityRealization} constructs an exact realization that keeps the terminals distinct.
\item \textbf{Progression.} \code{executedTrace} retains the stages; the obligation is carried to an arbitrary stage by the theorem identified above.
\end{enumerate}

\paragraph{Verification status.}
The definitions and proof bodies cited above were compared with the public
repository files at the revision stated on the first page. The sources
contain \texttt{\#print axioms} directives for the central declarations.
Their presence alone does not constitute an audit report of the results;
this manuscript does not claim to have independently executed
\texttt{lake build} or reproduced the outputs of the axiom audit.

\paragraph{Document status.}
This text reorganizes, without increasing the scope of the results, the
framework theorem
\Source{docs/theoreme-cadre-progression-godelienne-separation-exacte.fr.tex}{LaTeX}
and the canonical presentation
\href{https://github.com/JohnDoe-collab-stack/relational-perimeter/blob/1e5aeda745559d62a3819bb3213da7fce3c31ad1/docs/progression-constitutive-theories-godeliennes.fr.md}{\emph{Constitutive progression of theories and Gödelian obstructions}}.
Above all, it makes explicit the execution--contract--universal-obligation
axis that motivates the present order of exposition.

\section{Conclusion}

The shift established by the formalization starts from the execution
itself. A certified Gödelian production constitutes the objects, the
contract, and the finite witness that then make it possible to prove an
obligation governing every admissible exact realization, whatever its
memory type. The universal is therefore neither added externally to the
execution nor assumed in the definition of exactness: it is generated by
exact preservation of the distinctions that this particular execution
actually produced.

The finite progression repeats this passage without independently
reconstructing its inputs: every new stage consumes the retained history and
execution, constitutes its own contract, and then renews the universal
obligation. The central result therefore remains the one announced on the
first page:
\[
\boxed{\MainAxis}
\]

\end{document}
